\documentclass[../../main.tex]{subfiles}

\begin{document}

\section{测度保持变换与遍历性: Bogoliubov-Krilov 定理}\label{section:ch22-ergodic}

对可测空间 $(X,\mathcal{M})$, 映射 $T:X\to X$ 称为是可测变换, 若对每个可测集 $E$, $T^{-1}(E)$ 也是可测的. 观察到对映射 $T:X\to X$,
\begin{align}
T\text{ 是可测的当且仅当只要函数 }g\text{ 是可测的, }g\circ T\text{ 就可测}.
\label{align:ch22-measurable-transform}
\end{align}
对测度空间 $(X,\mathcal{M},\mu)$, 可测变换映射 $T:X\to X$ 称为是保测的, 若
\begin{align*}
\text{对所有 }A\in\mathcal{M},\quad\mu(T^{-1}(A))=\mu(A).
\end{align*}

\begin{proposition}[保测变换的积分刻画]\label{proposition:ch22-measure-preserving-integral}
令 $(X,\mathcal{M},\mu)$ 为有限测度空间而映射 $T:X\to X$ 是可测变换. 则 $T$ 是保测的当且仅当只要 $g$ 在 $X$ 上是可积的, $g\circ T$ 就在 $X$ 上可积, 且
\begin{align}
\text{对所有 }g\in L^{1}(X,\mu),\quad\int_{X}g\circ T\,d\mu=\int_{X}g\,d\mu.
\label{align:ch22-preserving-integral}
\end{align}
\end{proposition}

\begin{proof}
首先假设\eqref{align:ch22-preserving-integral}成立. 对 $A\in\mathcal{M}$, 由于 $\mu(X)<\infty$, 函数 $g=\chi_{A}$ 属于 $L^{1}(X,\mu)$ 且 $g\circ T=\chi_{T^{-1}(A)}$. 我们从\eqref{align:ch22-preserving-integral}推出 $\mu(T^{-1}(A))=\mu(A)$.

反过来, 假设 $T$ 保持测度. 令 $g$ 在 $X$ 上可积. 若 $g^{+}$ 是 $g$ 的正部, 则 $(g\circ T)^{+}=g^{+}\circ T$. 对负部, 类似的等式成立. 因此我们可以假设 $g$ 是非负的. 对简单函数 $g=\sum_{k=1}^{n}c_{k}\cdot\chi_{A_{k}}$, 由于 $T$ 是保持测度的,
\begin{align*}
\int_{X}g\circ T\,d\mu=\int_{X}\left[\sum_{k=1}^{n}c_{k}\cdot\chi_{A_{k}}\circ T\right]d\mu=\int_{X}\left[\sum_{k=1}^{n}c_{k}\cdot\chi_{T^{-1}(A_{k})}\right]d\mu\\
=\sum_{k=1}^{n}c_{k}\cdot\mu(A_{k})=\int_{X}g\,d\mu.
\end{align*}
因此对简单函数 $g$, \eqref{align:ch22-preserving-integral}成立. 根据简单逼近定理, 存在 $X$ 上的逐点收敛于 $g$ 的简单函数序列 $\{g_{n}\}$. 因此 $\{g_{n}\circ T\}$ 是 $X$ 上的逐点收敛于 $g\circ T$ 的简单函数序列. 用单调收敛定理两次和\eqref{align:ch22-preserving-integral}对简单函数成立, 有
\begin{align*}
\int_{X}g\circ T\,d\mu=\lim_{n\to\infty}\left[\int_{X}g_{n}\circ T\,d\mu\right]=\lim_{n\to\infty}\left[\int_{X}g_{n}\,d\mu\right]=\int_{X}g\,d\mu.
\end{align*}
\end{proof}

对测度空间 $(X,\mathcal{M},\mu)$ 和可测变换 $T:X\to X$, 可测集 $A$ 称为在 $T$ 下不变(关于 $\mu$), 若
\begin{align*}
\mu(A\sim T^{-1}(A))=\mu(T^{-1}(A)\sim A)=0,
\end{align*}
即模掉测度为零的集合, $T^{-1}(A)=A$. 显然,
\begin{align}
A\text{ 在 }T\text{ 下是不变的当且仅当在 }X\text{ 上 a.e. }\chi_{A}\circ T=\chi_{A}.
\label{align:ch22-invariant-set-char}
\end{align}
若 $(X,\mathcal{M},\mu)$ 也是一个概率空间, 即 $\mu(X)=1$, 变换 $T$ 称为是遍历的, 若对任何在 $T$ 下不变的集合 $A$, 有 $\mu(A)=0$ 或 $\mu(A)=1$.

\begin{proposition}[遍历性的函数刻画]\label{proposition:ch22-ergodic-char}
令 $(X,\mathcal{M},\mu)$ 为概率空间而 $T:X\to X$ 是保测变换. 则在 $X$ 上的实值可测函数 $g$ 中,
\begin{align}
T\text{ 是遍历的当且仅当只要在 }X\text{ 上 a.e. }g\circ T=g,\ g\text{ 就在 }X\text{ 上 a.e. 是常数}.
\label{align:ch22-ergodic-func}
\end{align}
\end{proposition}

\begin{proof}
首先假设只要在 $X$ 上 a.e. $g\circ T=g$, $g$ 就在 $X$ 上 a.e. 是常数. 令 $A\in\mathcal{M}$ 在 $T$ 下不变. 则 $g=\chi_{A}$($A$ 的特征函数)是可测的, 且在 $X$ 上 a.e. $\chi_{A}\circ T=\chi_{A}$. 因此 $\chi_{A}$ a.e. 是常数, 即 $\mu(A)=0$ 或 $\mu(A)=1$.

反过来, 假设 $T$ 是遍历的. 令 $g$ 为 $X$ 上的可测实值函数使得在 $X$ 上 a.e. $g\circ T=g$. 令 $k$ 为整数. 定义 $X_{k}=\{x\in X\mid k\leqslant g(x)<k+1\}$. 则 $X_{k}$ 是在 $T$ 下不变的可测集. 由 $T$ 的遍历性, 或者 $\mu(X_{k})=0$ 或者 $\mu(X_{k})=1$. 可数族 $\{X_{k}\}_{k\in\mathbf{Z}}$ 是不交的且它的并是 $X$. 由于 $\mu(X)=1$ 且 $\mu$ 是可数可加的, $\mu(X_{k})=0$. 定义 $I_{1}=[k',k'+1]$. 则 $\mu\{x\in X\mid g(x)\in I_{1}\}=1$ 且 $I_{1}$ 的长度 $\ell(I_{1})$是 $1$.

令 $n$ 为自然数使得下降的闭有界区间族 $\{I_{k}\}_{k=1}^{n}$ 对 $1\leqslant k\leqslant n$ 已定义, 满足
\begin{align*}
\ell(I_{k})=1/2^{k-1}\text{ 且 }\mu\{x\in X\mid g(x)\in I_{k}\}=1.
\end{align*}
令 $I_{n}=[a_{n},b_{n}]$, 定义 $c_{n}=(b_{n}-a_{n})/2$,
\begin{align*}
A_{n}=\{x\in X\mid a_{n}\leqslant g(x)<c_{n}\}\text{ 而 }B_{n}=\{x\in X\mid c_{n}\leqslant g(x)\leqslant b_{n}\}.
\end{align*}
则 $A_{n}$ 和 $B_{n}$ 是不交可测集, 其并为 $\{x\in X\mid g(x)\in I_{n}\}$, 即测度为 $1$ 的集合. 由于 $A_{n}$ 和 $B_{n}$ 在 $T$ 下是不变的, 我们从 $T$ 的遍历性推出这两个集合中恰有一个集合测度为 $1$. 若 $\mu(A_{n})=1$, 定义 $I_{n+1}=[a_{n},c_{n}]$. 否则, 定义 $I_{n+1}=[c_{n},b_{n}]$. 则 $\ell(I_{n+1})=1/2^{n}$ 且 $\mu\{x\in X\mid g(x)\in I_{n+1}\}=1$. 我们归纳地定义了闭有界区间的下降的可数族 $\{I_{n}\}_{n=1}^{\infty}$ 使得
\begin{align*}
\text{对所有 }n,\ \ell(I_{n})=1/2^{n-1}\text{ 且 }\mu\{x\in X\mid g(x)\in I_{n}\}=1.
\end{align*}
根据关于实数的集套定理, 存在属于每个 $I_{n}$ 的数 $c$. 我们宣称在 $X$ 上 $g=c$ a.e.. 事实上, 观察到若 $g(x)$ 属于 $I_{n}$, 则 $|g(x)-c|\leqslant 1/2^{n-1}$, 因此,
\begin{align*}
1=\mu\{x\in X\mid g(x)\in I_{n}\}\leqslant\mu\{x\in X\mid|g(x)-c|\leqslant 1/2^{n-1}\}\leqslant 1.
\end{align*}
由于
\begin{align*}
\{x\in X\mid g(x)=c\}=\bigcap_{n=1}^{\infty}\{x\in X\mid|g(x)-c|\leqslant 1/2^{n-1}\},
\end{align*}
我们从测度的连续性推出,
\begin{align*}
\mu\{x\in X\mid g(x)=c\}=\lim_{n\to\infty}\mu\{x\in X\mid|g(x)-c|\leqslant 1/2^{n-1}\}=1.
\end{align*}
\end{proof}

\begin{theorem}[Bogoliubov-Krilov]\label{theorem:ch22-bogoliubov-krilov}
令 $X$ 为紧度量空间而映射 $f:X\to X$ 是连续的. 则存在 Borel $\sigma$ 代数 $\mathcal{B}(X)$ 上的概率测度 $\mu$ 使得 $f$ 关于它是保测的.
\end{theorem}

\begin{proof}
考虑 $X$ 上具有最大值范数的连续实值函数的 Banach 空间 $C(X)$. 由于 $X$ 是紧度量空间, Borsuk 定理告诉我们 $C(X)$ 是可分的. 令 $\eta$ 为 $\mathcal{B}(X)$ 上的任何概率测度. 定义 $C(X)$ 上的线性泛函序列 $\{\psi_{n}\}$ 为
\begin{align}
\text{对所有 }n\in\mathbf{N}\text{ 和 }g\in C(X),\quad\psi_{n}(g)=\int_{X}\left[\frac{1}{n}\sum_{k=0}^{n-1}g\circ f^{k}\right]d\eta.
\label{align:ch22-cyclic-average}
\end{align}
观察到
\begin{align*}
\text{对所有 }n\in\mathbf{N}\text{ 和 }g\in C(X),\quad|\psi_{n}(g)|\leqslant\|g\|_{\max}.
\end{align*}
因此 $\{\psi_{n}\}$ 是 $[C(X)]^{*}$ 中的有界序列. 由于 Banach 空间 $C(X)$ 是可分的, 我们从 Helly 定理推出存在 $\{\psi_{n}\}$ 的子序列 $\{\psi_{n_{k}}\}$, 它关于弱 $*$ 拓扑收敛到有界泛函 $\psi\in[C(X)]^{*}$, 即
\begin{align*}
\text{对所有 }g\in C(X),\quad\lim_{k\to\infty}\psi_{n_{k}}(g)=\psi(g).
\end{align*}
因此对所有 $g\in C(X)$,
\begin{align*}
\lim_{k\to\infty}\psi_{n_{k}}(g\circ f)=\psi(g\circ f).
\end{align*}
然而, 对每个 $k$ 和 $g\in C(X)$,
\begin{align*}
\psi_{n_{k}}(g\circ f)-\psi_{n_{k}}(g)=\frac{1}{n_{k}}\left[\int_{X}[g\circ f^{n_{k}+1}-g]d\eta\right].
\end{align*}
取 $k\to\infty$ 的极限得到
\begin{align}
\text{对所有 }g\in C(X),\quad\psi(g\circ f)=\psi(g).
\label{align:ch22-invariant-psi}
\end{align}
由于每个 $\psi_{n}$ 是正泛函, 极限泛函 $\psi$ 也是正的. Riesz-Markov 定理告诉我们存在 Borel 测度 $\mu$ 使得
\begin{align*}
\text{对所有 }g\in C(X),\quad\psi(g)=\int_{X}g\,d\mu.
\end{align*}
我们从\eqref{align:ch22-invariant-psi}推出
\begin{align*}
\text{对所有 }g\in C(X),\quad\int_{X}g\circ f\,d\mu=\int_{X}g\,d\mu.
\end{align*}
根据命题\cref{proposition:ch22-measure-preserving-integral}, $f$ 关于 $\mu$ 是保测的. 最后, 对常函数 $g=1$, 对所有 $n$, $\psi_{n}(g)=1$, 因此 $\psi(g)=1$, 即 $\mu$ 是概率测度.
\end{proof}

\begin{proposition}[极值点与遍历性]\label{proposition:ch22-extreme-ergodic}
令 $f:X\to X$ 为紧度量空间 $X$ 上的连续映射. 定义 $\mathcal{M}_{f}$ 为 $\mathcal{B}(X)$ 上使得 $f$ 保测的概率测度集. 则 $\mathcal{M}_{f}$ 中的测度 $\mu$ 是 $\mathcal{M}_{f}$ 的极值点当且仅当 $f$ 关于 $\mu$ 是遍历的.
\end{proposition}

\begin{proof}
首先假设 $\mu$ 是 $\mathcal{M}_{f}$ 的极值点. 为证明 $f$ 是遍历的, 假设结论不成立. 则存在 $X$ 的 Borel 子集 $A$, 它在 $f$ 下关于 $\mu$ 是不变的, 而仍然有 $0<\mu(A)<1$. 对所有 $E\in\mathcal{B}(x)$ 定义
\begin{align*}
\nu(E)=\frac{\mu(E\cap A)}{\mu(A)}\text{ 与 }\eta(E)=\frac{\mu(E\cap[X\sim A])}{\mu(X\sim A)}.
\end{align*}
则由于 $\mu(X)=1$,
\begin{align*}
\mu=\lambda\nu+(1-\lambda)\eta,\text{ 其中 }\lambda=\mu(A).
\end{align*}
$\nu$ 和 $\eta$ 都是 $\mathcal{B}(X)$ 上的 Borel 概率测度. 我们宣称 $f$ 关于这些测度是保测的. 事实上, 由于 $f$ 关于测度 $\mu$ 保测而 $A$ 在 $f$ 下关于 $\mu$ 是不变的, 对每个 $E\in\mathcal{B}(X)$,
\begin{align*}
\mu(E\cap A)=\mu(f^{-1}(E\cap A))=\mu(f^{-1}(E)\cap f^{-1}(A))=\mu(f^{-1}(E)\cap A).
\end{align*}
因此 $f$ 关于 $\nu$ 是不变的. 用类似的方法可以证明, 它关于 $\eta$ 也是不变的. 因此 $\nu$ 和 $\eta$ 属于 $\mathcal{M}_{f}$, 且因此 $\mu$ 不是 $\mathcal{M}_{f}$ 的极值点. 因此 $f$ 是遍历的.

现在假定 $f$ 关于 $\mu\in\mathcal{M}_{f}$ 是遍历的. 为证明 $\mu$ 是 $\mathcal{M}_{f}$ 的极值点, 令 $\lambda\in(0,1)$ 而 $\nu,\eta\in\mathcal{M}_{f}$ 使得
\begin{align}
\mu=\lambda\nu+(1-\lambda)\eta.
\label{align:ch22-convex-decomp}
\end{align}
测度 $\nu$ 关于 $\mu$ 是绝对连续的. 由于 $\mu(X)<\infty$, Radon-Nikodym 定理告诉我们存在函数 $h\in L^{1}(X,\mu)$ 使得
\begin{align*}
\text{对所有 }A\in\mathcal{B}(X),\quad\nu(A)=\int_{A}h\,d\mu.
\end{align*}
从简单逼近定理和有界收敛定理得出,
\begin{align}
\text{对所有 }g\in L^{\infty}(X,\mu),\quad\int_{X}g\,d\nu=\int_{X}g\cdot h\,d\mu.
\label{align:ch22-rn-density}
\end{align}
固定 $\varepsilon>0$, 且定义 $X_{\varepsilon}=\{x\in X\mid h(x)\geqslant 1/\lambda+\varepsilon\}$. 我们从\eqref{align:ch22-convex-decomp}推出,
\begin{align*}
\mu(X_{\varepsilon})\geqslant\lambda\cdot\int_{X_{\varepsilon}}h\,d\mu\geqslant(1+\lambda\cdot\varepsilon)\cdot\mu(X_{\varepsilon}).
\end{align*}
因此 $\mu(X_{\varepsilon})=0$. $h$ 和 $h\circ f$ 在 $X$ 上关于 $\mu$ 是本性有界的. 因此, 用\eqref{align:ch22-rn-density}, 首先对 $g=h\circ f$ 接着对 $g=h$, 以及 $f$ 关于 $\nu$ 的不变性, 有
\begin{align*}
\int_{X}h\circ f\cdot h\,d\mu=\int_{X}h\circ f\,d\nu=\int_{X}h\,d\nu=\int_{X}h^{2}\,d\mu.
\end{align*}
我们从该等式和 $f$ 关于 $\mu$ 的不变性推出,
\begin{align*}
\int_{X}[h\circ f-h]^{2}\,d\mu=\int_{X}[h\circ f]^{2}\,d\mu-2\cdot\int_{X}h\circ f\cdot h\,d\mu+\int_{X}h^{2}\,d\mu\\
=2\cdot\int_{X}h^{2}\,d\mu-2\cdot\int_{X}h\circ f\cdot h\,d\mu\\
=2\cdot\int_{X}h^{2}\,d\mu-2\cdot\int_{X}h^{2}\,d\mu=0.
\end{align*}
因此在 $X$ 上 $h\circ f=h$ a.e. $[\mu]$. 根据 $f$ 的遍历性和命题\cref{proposition:ch22-ergodic-char}, 存在常数 $c$ 使得在 $X$ 上 $h=c$ a.e. $[\mu]$. 但 $\mu$ 和 $\nu$ 是概率测度, 因此,
\begin{align*}
1=\nu(X)=\int_{X}h\,d\mu=c\cdot\mu(X)=c.
\end{align*}
因而 $\mu=\nu$, 于是 $\mu=\eta$. 因此 $\mu$ 是 $\mathcal{M}_{f}$ 的极值点.
\end{proof}

\begin{theorem}[遍历概率测度的存在性]\label{theorem:ch22-ergodic-existence}
令 $f:X\to X$ 为紧度量空间 $X$ 上的连续映射. 则存在 Borel $\sigma$ 代数 $\mathcal{B}(X)$ 上的概率测度 $\mu$ 使得 $f$ 在其上是遍历的.
\end{theorem}

\begin{proof}
令 $\mathcal{R}adon(X)$ 为 $\mathcal{B}(X)$ 上的带号 Radon 测度的 Banach 空间, 而线性算子 $\Phi:\mathcal{R}adon(X)\to[C(X)]^{*}$ 定义为
\begin{align*}
\text{对所有 }\mu\in\mathcal{R}adon(X)\text{ 与 }g\in C(X),\quad\Phi(\mu)(g)=\int_{X}g\,d\mu.
\end{align*}
关于 $C(X)$ 的对偶的 Riesz 表示定理告诉我们, $\Phi$ 是 $\mathcal{R}adon(X)$ 映上 $[C(X)]^{*}$ 的线性同构. 定义 $\mathcal{M}_{f}$ 为 $\mathcal{B}(X)$ 上使得 $f$ 保测的概率测度集. 则测度 $\mu$ 是 $\mathcal{M}_{f}$ 的极值点当且仅当 $\Phi(\mu)$ 是 $\Phi(\mathcal{M}_{f})$ 的极值点. 因此, 根据前一个命题, 为证明定理我们必须证明集合 $\Phi(\mathcal{M}_{f})$ 具有极值点. 根据 Bogoliubov-Krilov 定理, $\mathcal{M}_{f}$ 是非空的. Krein-Milman 定理的一个推论, 即前一章的推论 $13$, 告诉我们若 $\Phi(\mathcal{M}_{f})$ 是有界的、凸的, 且关于弱 $*$ 拓扑是闭的, 则 $\mathcal{M}_{f}$ 具有极值点. Riesz-Markov 定理告诉我们 $\Phi$ 定义了 Radon 测度映上正泛函的同构. 正泛函当然是弱 $*$ 闭的, 在常函数 $1$ 上取值 $1$ 的泛函也如此. 根据命题\cref{proposition:ch22-ergodic-char}, 泛函 $\psi\in[C(X)]^{*}$ 是关于 $f$ 不变的测度在 $\Phi$ 下的象当且仅当
\begin{align*}
\text{对所有 }g\in C(X),\quad\psi(g\circ f)-\psi(g)=0.
\end{align*}
固定 $g\in C(X)$. 函数 $g\circ f-g$ 的估值是 $[C(X)]^{*}$ 上关于弱 $*$ 拓扑连续的线性泛函, 因此它的核是弱 $*$ 闭的. 因此交
\begin{align*}
\bigcap_{g\in C(X)}\{\psi\in[C(X)]^{*}\mid\psi(g\circ f)=\psi(g)\}
\end{align*}
也是弱 $*$ 闭集. 这完成了 $\Phi(\mathcal{M}_{f})$ 的弱 $*$ 闭性的证明, 因此完成了定理的证明.
\end{proof}

渐近平均现象最初在气体的动态分析中被引入. 以下定理揭示在这种现象的研究中遍历性的意义. 观察到\eqref{align:ch22-birkhoff}的右边与点 $x\in X$ 无关.

\begin{theorem}[Birkhoff 遍历定理]\label{theorem:ch22-birkhoff}
令 $T$ 为概率空间 $(X,\mathcal{M},\mu)$ 上的测度保持变换. 则 $T$ 是遍历的当且仅当对每个 $g\in L^{1}(X,\mu)$, 对几乎所有 $x\in X$,
\begin{align}
\lim_{n\to\infty}\left[\frac{1}{n}\sum_{k=0}^{n-1}g(T^{k}(x))\right]=\frac{1}{\mu(X)}\int_{X}g\,d\mu.
\label{align:ch22-birkhoff}
\end{align}
该定理的证明可在 Michael Brin 和 Garrett Stuck 的书《Introduction to Dynamical Systems》[\cite{BS02}]与 Paul Halmos 的《Lectures on Ergodic Theory》[\cite{Hal06}]中找到. 这些书也包含了测度保持和遍历变换的多个例子.
\end{theorem}

\begin{exercise}\label{exercise:ch22-26}
Bogoliubov-Krilov 定理的证明是否也提供了 $X$ 是紧 Hausdorff 空间但不必是可度量化的情形的证明?
\end{exercise}

\begin{exercise}\label{exercise:ch22-27}
令 $(X,\mathcal{M},\mu)$ 为有限测度空间而 $T:X\to X$ 是可测变换. 对 $X$ 上的可测函数 $g$, 定义可测函数 $U_{T}(g)$ 为 $U_{T}(g)(x)=g(T(x))$. 证明 $T$ 是测度保持的当且仅当对每个 $1\leqslant p<\infty$, $U_{T}$ 将 $L^{p}(X,\mu)$ 映到其自身且 $U_{T}$ 是等距的.
\end{exercise}

\begin{exercise}\label{exercise:ch22-28}
假定 $T:\mathbf{R}^{n}\to\mathbf{R}^{n}$ 是线性的. 对 $T$ 建立使得它关于 $\mathbf{R}^{n}$ 上的 Lebesgue 测度保测的必要与充分条件.
\end{exercise}

\begin{exercise}\label{exercise:ch22-29}
令 $S^{1}=\{z=e^{i\theta}\mid\theta\in\mathbf{R}\}$ 为以复数的乘法为群运算的圆周, 而 $\mu$ 为该群上的 Haar 测度(见习题\cref{exercise:ch22-24}). 定义 $T:S^{1}\to S^{1}$ 为 $T(z)=z^{2}$. 证明 $T$ 保持 $\mu$.
\end{exercise}

\begin{exercise}\label{exercise:ch22-30}
定义 $f:\mathbf{R}^{2}\to\mathbf{R}^{2}$ 为 $f(x,y)=(2x,y/2)$. 证明 $f$ 关于 Lebesgue 测度保测.
\end{exercise}

\begin{exercise}\label{exercise:ch22-31}
(Poincar\'e 回复)令 $T$ 为有限测度空间 $(X,\mathcal{M},\mu)$ 上的测度保持变换, 而集合 $A$ 是可测的. 证明对几乎所有 $x\in A$, 存在无穷多个自然数 $n$ 使得 $T^{n}(x)$ 属于 $A$.
\end{exercise}

\end{document}